% GENERATED FILE. Edit canonical lessons, then run npm run generate:books.
% canonical-source-hash: fnv1a64:d1236c6eee843b87
% canonical-lessons: KA-C269-varadi, KA-C269-samasye, KA-C269-suddi, KA-C269-mahiti, KA-C269-anka

\chapter{Report, Problem, News, Information, Mark}
\label{ch:ka-report-problem-news-information-mark}
\hlchaptermodality{269}

\begin{chapteropening}
\textbf{By the end of this chapter:} \emph{I can say a report, a problem, news, information and a mark.}

Complete the last lesson of chapter 269: I can say a report, a problem, news, information and a mark.
\end{chapteropening}

\section[\texorpdfstring{\kn{ವರದಿ} (varadi)}{ವರದಿ (varadi)}]{\kn{ವರದಿ} (varadi) — a report}

\label{lesson:KA-C269-varadi}



\begin{quote}
Before the new one: say the Kannada for a company, then the Kannada for a class.
\end{quote}

\subsection*{You'll want to know: \kn{ವರದಿ}}
\textbf{\kn{ವರದಿ}} — \emph{varadi} — \textquotedblleft{}a report\textquotedblright{}.

\begin{grammarlens}[title={What you've built}]
One more word for this chapter.
\end{grammarlens}

\subsection*{Your turn}
\begin{itemize}
\raggedright
  \item \emph{Say it:} \emph{varadi}
  \item \emph{Say it:} \emph{varadi}, once more
  \item \emph{Say it:} say \emph{varadi}
  \item \emph{Recall:} say the Kannada for a company, then the Kannada for a class, then say \emph{varadi} again
\end{itemize}

\begin{tcolorbox}[breakable,colback=teal!4,colframe=teal!35!black,title={Before you move on}]
What does \kn{ವರದಿ} mean? (\textquotedblleft{}A report\textquotedblright{}.) Say it once more.
\end{tcolorbox}
\section[\texorpdfstring{\kn{ಸಮಸ್ಯೆ} (samasye)}{ಸಮಸ್ಯೆ (samasye)}]{\kn{ಸಮಸ್ಯೆ} (samasye) — a problem}

\label{lesson:KA-C269-samasye}



\begin{quote}
Before the new one: say the Kannada for a class, then the Kannada for a report.
\end{quote}

\subsection*{You'll want to know: \kn{ಸಮಸ್ಯೆ}}
\textbf{\kn{ಸಮಸ್ಯೆ}} — \emph{samasye} — \textquotedblleft{}a problem\textquotedblright{}.

\begin{grammarlens}[title={What you've built}]
One more word for this chapter.
\end{grammarlens}

\subsection*{Your turn}
\begin{itemize}
\raggedright
  \item \emph{Say it:} \emph{samasye}
  \item \emph{Say it:} \emph{samasye}, once more
  \item \emph{Say it:} say \emph{samasye}
  \item \emph{Recall:} say the Kannada for a class, then the Kannada for a report, then say \emph{samasye} again
\end{itemize}

\begin{tcolorbox}[breakable,colback=teal!4,colframe=teal!35!black,title={Before you move on}]
What does \kn{ಸಮಸ್ಯೆ} mean? (\textquotedblleft{}A problem\textquotedblright{}.) Say it once more.
\end{tcolorbox}
\section[\texorpdfstring{\kn{ಸುದ್ದಿ} (suddi)}{ಸುದ್ದಿ (suddi)}]{\kn{ಸುದ್ದಿ} (suddi) — news}

\label{lesson:KA-C269-suddi}



\begin{quote}
Before the new one: say the Kannada for a report, then the Kannada for a problem.
\end{quote}

\subsection*{You'll want to know: \kn{ಸುದ್ದಿ}}
\textbf{\kn{ಸುದ್ದಿ}} — \emph{suddi} — \textquotedblleft{}news\textquotedblright{}.

\begin{grammarlens}[title={What you've built}]
One more word for this chapter.
\end{grammarlens}

\subsection*{Your turn}
\begin{itemize}
\raggedright
  \item \emph{Say it:} \emph{suddi}
  \item \emph{Say it:} \emph{suddi}, once more
  \item \emph{Say it:} say \emph{suddi}
  \item \emph{Recall:} say the Kannada for a report, then the Kannada for a problem, then say \emph{suddi} again
\end{itemize}

\begin{tcolorbox}[breakable,colback=teal!4,colframe=teal!35!black,title={Before you move on}]
What does \kn{ಸುದ್ದಿ} mean? (\textquotedblleft{}News\textquotedblright{}.) Say it once more.
\end{tcolorbox}
\section[\texorpdfstring{\kn{ಮಾಹಿತಿ} (māhiti)}{ಮಾಹಿತಿ (māhiti)}]{\kn{ಮಾಹಿತಿ} (māhiti) — information}

\label{lesson:KA-C269-mahiti}



\begin{quote}
Before the new one: say the Kannada for a problem, then the Kannada for news.
\end{quote}

\subsection*{You'll want to know: \kn{ಮಾಹಿತಿ}}
\textbf{\kn{ಮಾಹಿತಿ}} — \emph{māhiti} — \textquotedblleft{}information\textquotedblright{}.

\begin{grammarlens}[title={What you've built}]
One more word for this chapter.
\end{grammarlens}

\subsection*{Your turn}
\begin{itemize}
\raggedright
  \item \emph{Say it:} \emph{māhiti}
  \item \emph{Say it:} \emph{māhiti}, once more
  \item \emph{Say it:} say \emph{māhiti}
  \item \emph{Recall:} say the Kannada for a problem, then the Kannada for news, then say \emph{māhiti} again
\end{itemize}

\begin{tcolorbox}[breakable,colback=teal!4,colframe=teal!35!black,title={Before you move on}]
What does \kn{ಮಾಹಿತಿ} mean? (\textquotedblleft{}Information\textquotedblright{}.) Say it once more.
\end{tcolorbox}
\section[\texorpdfstring{\kn{ಅಂಕ} (aṅka)}{ಅಂಕ (aṅka)}]{\kn{ಅಂಕ} (aṅka) — a mark, a score}

\label{lesson:KA-C269-anka}



\begin{quote}
Before the new one: say the Kannada for news, then the Kannada for information.
\end{quote}

\subsection*{You'll want to know: \kn{ಅಂಕ}}
\textbf{\kn{ಅಂಕ}} — \emph{aṅka} — \textquotedblleft{}a mark, a score\textquotedblright{}.

\begin{grammarlens}[title={What you've built}]
One more word for this chapter.
\end{grammarlens}

\subsection*{Your turn}
\begin{itemize}
\raggedright
  \item \emph{Say it:} \emph{aṅka}
  \item \emph{Say it:} \emph{aṅka}, once more
  \item \emph{Say it:} say \emph{aṅka}
  \item \emph{Recall:} say the Kannada for news, then the Kannada for information, then say \emph{aṅka} again
\end{itemize}

\begin{tcolorbox}[breakable,colback=teal!4,colframe=teal!35!black,title={Before you move on}]
What does \kn{ಅಂಕ} mean? (\textquotedblleft{}A mark, a score\textquotedblright{}.) Say it once more.
\end{tcolorbox}
